\labsection{7}{Tree representation and traversal}{sec:lab07}

\begin{sourcealignment}{Official Lab Manual --- Lab VII}
\textbf{Objective.} Learn the tree ADT and apply trees to real-life problems.

\textbf{Outcomes.} Implement a tree using an array and a linked representation; implement traversal methods; choose an appropriate tree structure.

\textbf{Problem directions.} Tree representation; binary-tree implementation; traversal algorithms; inserting, deleting, and updating tree elements.
\end{sourcealignment}

\begin{labproject}{Represent the same binary tree two ways}
\textbf{Coverage.} Chapter 7.\\
\textbf{Goal.} Represent one binary tree in both linked and array form and produce all three traversal orders.

\begin{lstlisting}
TreeNode* root = new TreeNode{8, nullptr, nullptr};
root->left = new TreeNode{4, nullptr, nullptr};
root->right = new TreeNode{12, nullptr, nullptr};
root->left->left = new TreeNode{2, nullptr, nullptr};
root->left->right = new TreeNode{6, nullptr, nullptr};
\end{lstlisting}

Implement inorder, preorder, and postorder traversal. Then represent the same shape in an array and identify each parent/child index.

\begin{tryit}
Write all three traversal sequences before running the program. Add one update operation and one structural insertion/deletion exercise specified by your lab session.
\end{tryit}
\end{labproject}

\begin{verification}
\begin{itemize}
  \item Traversals visit every stored node exactly once.
  \item The linked and array representations describe the same shape.
  \item Your array child-index formula matches the chosen 0-based or 1-based convention.
  \item Empty subtrees terminate recursion safely.
\end{itemize}
\end{verification}

\begin{labdeliverable}
Submit both representations, all three traversal outputs, and a brief note stating your indexing convention.
\end{labdeliverable}
